% !TeX root = ../main.tex
% ABSTRACT: traducción del resumen. Se compone en inglés (guiones,
% comillas) y con punto decimal, también en \num y \qty de siunitx.
\begin{resumenproyecto}{ABSTRACT}
\begin{otherlanguage}{english}
\sisetup{output-decimal-marker={.}}

\apartadoresumen{1. Introduction}
Problem context, motivation and objectives.

\apartadoresumen{2. Methodology}
What was done and how.

\apartadoresumen{3. Results}
Key quantitative results, e.g.\ an error of \qty{3.9}{\percent}.

\apartadoresumen{4. Conclusions}
Main conclusions and open lines of work.

\vspace{1em}
\noindent\textbf{Keywords:} \Keywords

\end{otherlanguage}
\end{resumenproyecto}
